\documentclass[12pt]{article}

\usepackage[margin=1in]{geometry}
\usepackage{amsmath,amsthm,amssymb}
\usepackage{mathrsfs}
\usepackage{mathtools}
\usepackage{enumitem}
\usepackage{physics}

\newcommand{\magsq}[1]{\big|#1\big|^2}
\newcommand{\avg}[1]{\left<#1\right>}
\newcommand{\fullint}{\int_{-\infty}^\infty}
\newcommand{\fullintd}[1]{\fullint\dd#1\:}
\newcommand{\cint}[2]{\int_{#1}^{#2}}
\newcommand{\cintd}[3]{\cint{#1}{#2}\dd#3\:}
\newcommand\treq{\stackrel{\mathclap{\tiny\mbox{Tr}}}{=}}

\begin{document}
	
\title{Exercise Set 10}
\author{Sean Ericson \\ Phys 633}
\maketitle

\section*{Tuesday}

\subsection*{Exercise 1}
The initial joint state is given by
\[ \left(\frac{\ket{1_L0_R} + \ket{0_L1_R}}{\sqrt{2}}\right)\ket{\uparrow_A\uparrow_B} = \frac{\ket{1_L0_R}\ket{\uparrow_A\uparrow_B} + \ket{0_L1_R}\ket{\uparrow_A\uparrow_B}}{\sqrt{2}}. \]
After the interaction (and discarding the mode states) we have the state
\[ \frac{\ket{\downarrow_A\uparrow_B} + \ket{\uparrow_A\downarrow_B}}{\sqrt{2}}, \]
in which the spins of the particals are clearly entangled. Assuming the modes are physically separted, and the spin-flipping interactions between the modes and the particals are local, the entanglment present in the particals' spins must have come from the modes. Hence, mode state
\[ \frac{\ket{1_L0_R} + \ket{0_L1_R}}{\sqrt{2}} \]
must be entangled.


\end{document}